% ch27_weighted_limits.tex — Volume III Chapter 3

\chapter{Weighted Limits and Colimits}

\begin{overviewbox}
Chapter~20 introduced weighted limits in the enriched setting: a $W$-weighted
limit $\{W, F\}$ of $F : \mathcal{J} \to \mathcal{C}$ weighted by
$W : \mathcal{J} \to \mathcal{V}$ represents the $\mathcal{V}$-functor
$c \mapsto [\mathcal{J}, \mathcal{V}]_\mathcal{V}(W, \mathcal{C}(c, F-))$.
Ordinary limits were the case $W = \Delta I$.

This chapter uses the coend calculus of Chapter~26 to develop the full
theory. We re-derive weighted limits and colimits via ends and coends, prove
the central representability theorem, work out the principal examples
(cotensor, conical, lax limit, comma object), and show how all of ordinary
limits-and-colimits theory transports to the weighted setting. The chapter
closes with a preview of weighted Kan extensions (Chapter~31): in the enriched
world, \emph{every} (co)limit is a Kan extension and \emph{every} Kan
extension is a weighted (co)limit.
\end{overviewbox}


% ═══════════════════════════════════════════════════════════════════════════
\begin{nameorigin}
\term{Weighted limit} (also \emph{indexed limit}) generalizes ordinary
limits to the enriched setting by allowing the ``shape'' of the limit
to be modulated by a weighting functor $W : \mathcal{J} \to \mathcal{V}$.
The terminology is Kelly's (Vol II Ch.~20). The intuition: ordinary
limits compute by taking the ``equal-weight average'' of a diagram
(weight $W = $ const $\mathbf{1}$); weighted limits allow nontrivial
weight distributions. Modern enriched category theory works almost
entirely in the weighted setting, since ordinary limits are too
restrictive in $\mathcal{V}$-categories.
\end{nameorigin}

\section{Weighted Limits via Ends}
\label{sec:wlim-end}
% ═══════════════════════════════════════════════════════════════════════════

\subsection{Formalizing}

Fix a complete symmetric monoidal closed category $\mathcal{V}$ and a small
$\mathcal{V}$-category $\mathcal{J}$. Let $F : \mathcal{J} \to \mathcal{C}$
be a $\mathcal{V}$-functor (a \emph{diagram}) and $W : \mathcal{J} \to
\mathcal{V}$ another $\mathcal{V}$-functor (a \emph{weight}).

\begin{definition}[Weighted limit]
\label{def:weighted-limit}
The \term{weighted limit} of $F$ by $W$, written $\{W, F\}$ or $\lim^W F$,
is an object of $\mathcal{C}$ — when it exists — together with a $\mathcal{V}$-natural
isomorphism
\begin{equation*}
  \mathcal{C}(c, \{W, F\}) \;\cong\; [\mathcal{J}, \mathcal{V}]_\mathcal{V}(W,
  \mathcal{C}(c, F-))
\end{equation*}
in $\mathcal{V}$, naturally in $c \in \mathcal{C}$.
\end{definition}

\begin{howtosay}
$\{W, F\}$ \sayit{``the W-weighted limit of F''} or \sayit{``the limit of F
weighted by W''}.
\end{howtosay}

\begin{theorem}[End formula for weighted limits]
\label{thm:wlim-end-formula}
The right-hand side of Definition~\ref{def:weighted-limit} is an end:
\begin{equation*}
  [\mathcal{J}, \mathcal{V}]_\mathcal{V}(W, \mathcal{C}(c, F-))
  \;=\;
  \int_{j \in \mathcal{J}} \mathcal{V}(W(j), \mathcal{C}(c, F j)).
\end{equation*}
\end{theorem}

\begin{proof}
This is the end formula for $\mathcal{V}$-natural transformations
(Chapter~26, Proposition on enriched ends). Apply it with $G = W$ and $H =
\mathcal{C}(c, F-)$.
\end{proof}

Combining:
\begin{equation*}
  \mathcal{C}(c, \{W, F\})
  \;\cong\;
  \int_j \mathcal{V}(W j, \mathcal{C}(c, F j)).
\end{equation*}

\subsection{Reorganizing: Conical Limits as $\Delta I$-Weighted}

\begin{proposition}[Ordinary limits as weighted limits]
\label{prop:conical-as-weighted}
For $\mathcal{V} = \mathbf{Set}$ and $W = \Delta\{*\}$ the terminal weight,
$\{\Delta\{*\}, F\} = \lim F$, the ordinary (\term{conical}) limit.
\end{proposition}

\begin{proof}
The end formula gives
\begin{equation*}
  \mathcal{C}(c, \{\Delta\{*\}, F\})
  = \int_j \mathbf{Set}(\{*\}, \mathcal{C}(c, F j))
  = \int_j \mathcal{C}(c, F j)
  = \mathrm{Nat}(\Delta c, F)
  = \mathrm{cones}(c, F),
\end{equation*}
the set of cones from $c$ to $F$. By the universal property of the limit,
this is $\mathcal{C}(c, \lim F)$, so $\{\Delta\{*\}, F\} \cong \lim F$.
\end{proof}

\subsection{Formally: Other Weighted Limits Worth Naming}

\begin{example_thm}[Cotensor]
\label{ex:cotensor}
For a single-object indexing category $\mathcal{J} = \mathbf{1}$ (terminal
$\mathcal{V}$-category), $F$ is a single object $X \in \mathcal{C}$, $W$ is a
single object $V \in \mathcal{V}$, and the weighted limit is the
\term{cotensor} $V \pitchfork X$ (also written $X^V$ or $[V, X]_\mathcal{C}$):
\begin{equation*}
  \mathcal{C}(c, V \pitchfork X) \;\cong\; \mathcal{V}(V, \mathcal{C}(c, X)).
\end{equation*}
When $\mathcal{C} = \mathcal{V}$ self-enriched, this is the internal hom
$[V, X]$.
\end{example_thm}

\begin{example_thm}[Lax limit]
For $\mathcal{J} = \mathbf{2}$ (the walking arrow) and $W : \mathbf{2} \to
\mathbf{Cat}$ sending the arrow to a specific 1-cell, the weighted limit is
the \term{lax limit} or \emph{comma object}. In $\mathbf{Cat}$, this is the
comma category $(F \downarrow G)$ when $F = W(0)$ and $G = W(1)$.
\end{example_thm}

\begin{example_thm}[Inserter, equifier, descent]
More elaborate weights generate \emph{inserters} (a 1-cell making two given
2-cells equal up to a new mediating 2-cell), \emph{equifiers} (objects that
equify two 2-cells), and \emph{descent objects}. These are the standard
\emph{flexible limits} of 2-category theory and arise as weighted limits
with explicit small $\mathbf{Cat}$-valued weights.
\end{example_thm}


% ═══════════════════════════════════════════════════════════════════════════
\section{Weighted Colimits via Coends}
\label{sec:wcolim-coend}
% ═══════════════════════════════════════════════════════════════════════════

\subsection{Formalizing}

\begin{definition}[Weighted colimit]
\label{def:weighted-colimit}
For a diagram $F : \mathcal{J} \to \mathcal{C}$ and weight $W :
\mathcal{J}^{\mathrm{op}} \to \mathcal{V}$ (note the variance), the
\term{weighted colimit} is an object $W \star F \in \mathcal{C}$ — when it
exists — with a $\mathcal{V}$-natural isomorphism
\begin{equation*}
  \mathcal{C}(W \star F, c) \;\cong\; [\mathcal{J}^{\mathrm{op}}, \mathcal{V}]_\mathcal{V}(W,
  \mathcal{C}(F-, c)).
\end{equation*}
\end{definition}

\begin{howtosay}
$W \star F$ \sayit{``W-weighted colimit of F''}. The star is the standard
shorthand for the weighted colimit construction.
\end{howtosay}

\begin{theorem}[Coend formula for weighted colimits]
\label{thm:wcolim-coend-formula}
In a tensored $\mathcal{V}$-category $\mathcal{C}$ (where $V \otimes c$ exists
for $V \in \mathcal{V}$ and $c \in \mathcal{C}$ with the adjunction
$\mathcal{C}(V \otimes c, d) \cong \mathcal{V}(V, \mathcal{C}(c, d))$), the
weighted colimit is given by the coend
\begin{equation*}
  W \star F \;\cong\; \int^{j} W(j) \otimes F(j).
\end{equation*}
\end{theorem}

\begin{prooflog}
\proofstep{Compute $\mathcal{C}(W \star F, c)$ using the proposed coend
formula.}{Substitute the candidate into the hom-functor.}
\proofstep{$\mathcal{C}\!\left(\int^j W(j) \otimes F(j), c\right) \cong
\int_j \mathcal{C}(W(j) \otimes F(j), c)$.}{Hom-functor turns coends into
ends in the first variable (continuity).}
\proofstep{$\int_j \mathcal{C}(W(j) \otimes F(j), c) \cong
\int_j \mathcal{V}(W(j), \mathcal{C}(F(j), c))$.}{By the tensor-cotensor
adjunction of Definition assumption.}
\proofstep{$\int_j \mathcal{V}(W(j), \mathcal{C}(F(j), c)) =
[\mathcal{J}^{\mathrm{op}}, \mathcal{V}]_\mathcal{V}(W, \mathcal{C}(F-, c))$.}{
End formula for enriched natural transformations (Theorem~\ref{thm:wlim-end-formula}).}
\proofstep{Therefore $\mathcal{C}\!\left(\int^j W(j) \otimes F(j), c\right)
\cong \mathcal{C}(W \star F, c)$ $\mathcal{V}$-naturally in $c$.}{Composing
the chain of isomorphisms.}
\proofstep{By the Yoneda lemma, $\int^j W(j) \otimes F(j) \cong W \star F$.}{
Two objects with $\mathcal{V}$-naturally isomorphic hom-out functors are
isomorphic.}
\end{prooflog}

\begin{corollary}[Conical colimit as weighted]
$\Delta I \star F \cong \operatorname{colim} F$ (the ordinary colimit), where
$\Delta I : \mathcal{J}^{\mathrm{op}} \to \mathcal{V}$ is the constant unit.
\end{corollary}

\subsection{Reorganizing: Tensor as Weighted Colimit}

\begin{example_thm}[Tensor]
For $\mathcal{J} = \mathbf{1}$, $F$ a single object $X$, $W$ a single object
$V$, the weighted colimit is the \term{tensor} $V \otimes X$ (also written
$V \cdot X$ in $\mathbf{Set}$-enriched settings):
\begin{equation*}
  \mathcal{C}(V \otimes X, c) \;\cong\; \mathcal{V}(V, \mathcal{C}(X, c)).
\end{equation*}
This is the dual of cotensor (Example~\ref{ex:cotensor}); together they
characterize $\mathcal{V}$-tensored-and-cotensored categories.
\end{example_thm}

\begin{example_thm}[Coend as weighted colimit of a profunctor]
For $T : \mathcal{C}^{\mathrm{op}} \otimes \mathcal{C} \to \mathcal{V}$, treating
$T$ as a diagram $\mathcal{C}^{\mathrm{op}} \otimes \mathcal{C} \to \mathcal{V}$
and the hom-profunctor $\mathcal{C}(-, -)$ as the weight, the coend $\int^c
T(c, c)$ is the weighted colimit of $T$ along the hom-weight.
\end{example_thm}


% ═══════════════════════════════════════════════════════════════════════════
\section{Functoriality and Yoneda Reduction}
\label{sec:wlim-functoriality}
% ═══════════════════════════════════════════════════════════════════════════

\subsection{Formalizing}

\begin{proposition}[Functoriality of weighted limits]
\label{prop:wlim-functorial}
$\{-, -\} : [\mathcal{J}, \mathcal{V}]^{\mathrm{op}} \otimes
[\mathcal{J}, \mathcal{C}] \to \mathcal{C}$ is a $\mathcal{V}$-functor:
weighted limits are functorial in both weight and diagram, contravariantly
in the weight.
\end{proposition}

\begin{proof}
The defining isomorphism $\mathcal{C}(c, \{W, F\}) \cong \int_j \mathcal{V}(Wj,
\mathcal{C}(c, Fj))$ is $\mathcal{V}$-natural in $c$, but also in $W$
(contravariantly) and $F$ (covariantly) by the universal property of the end.
Apply Yoneda in $c$ to lift naturality from the hom-functors to $\{W, F\}$
itself.
\end{proof}

\begin{theorem}[Yoneda reduction for weighted (co)limits]
\label{thm:wlim-yoneda-reduction}
For $F : \mathcal{J} \to \mathcal{C}$ and $j_0 \in \mathcal{J}$:
\begin{align*}
  \{\mathcal{J}(j_0, -), F\} &\;\cong\; F(j_0) \quad\text{(weighted limit)} \\
  \mathcal{J}(-, j_0) \star F &\;\cong\; F(j_0) \quad\text{(weighted colimit)}
\end{align*}
\end{theorem}

\begin{prooflog}
\proofstep{For the weighted limit: $\mathcal{C}(c, \{\mathcal{J}(j_0, -), F\})
\cong \int_j \mathcal{V}(\mathcal{J}(j_0, j), \mathcal{C}(c, Fj))$.}{Apply the
defining formula.}
\proofstep{$\int_j \mathcal{V}(\mathcal{J}(j_0, j), \mathcal{C}(c, Fj))$ is the
$\mathcal{V}$-natural transformations from $\mathcal{J}(j_0, -)$ to
$\mathcal{C}(c, F-)$.}{By the end formula for $\mathcal{V}$-naturals.}
\proofstep{By the (enriched) Yoneda lemma, this equals $\mathcal{C}(c, F(j_0))$.}{
Yoneda applied to the $\mathcal{V}$-functor $\mathcal{C}(c, F-) : \mathcal{J} \to
\mathcal{V}$.}
\proofstep{So $\mathcal{C}(c, \{\mathcal{J}(j_0, -), F\}) \cong \mathcal{C}(c, F(j_0))$
$\mathcal{V}$-naturally in $c$.}{Combining the above.}
\proofstep{By Yoneda in $c$, $\{\mathcal{J}(j_0, -), F\} \cong F(j_0)$.}{
Yoneda for objects.}
\proofstep{For the weighted colimit: $\mathcal{C}(\mathcal{J}(-, j_0) \star F, c)
\cong \int_j \mathcal{V}(\mathcal{J}(j, j_0), \mathcal{C}(F j, c))$.}{Apply
the defining formula.}
\proofstep{This is the $\mathcal{V}$-naturals $\mathcal{J}(-, j_0) \Rightarrow
\mathcal{C}(F-, c)$.}{End formula again.}
\proofstep{By contravariant Yoneda, equals $\mathcal{C}(F(j_0), c)$.}{Yoneda
applied to $\mathcal{C}(F-, c) : \mathcal{J}^{\mathrm{op}} \to \mathcal{V}$.}
\proofstep{Hence $\mathcal{J}(-, j_0) \star F \cong F(j_0)$.}{Yoneda for
objects.}
\end{prooflog}

\begin{remark}[Representable weights pick out values]
Theorem~\ref{thm:wlim-yoneda-reduction} says: \emph{representable weights
evaluate the diagram}. This is the weighted analogue of the fact that
``the limit over a singleton diagram is the object itself.''
\end{remark}

\subsection{Reorganizing: Bilinearity and Distributivity}

\begin{proposition}[Colimits-into-limits commutation]
\label{prop:bilinearity}
For $W : \mathcal{J} \to \mathcal{V}$, $F : \mathcal{J} \to \mathcal{C}$, and
$G : \mathcal{K} \to \mathcal{C}$:
\begin{equation*}
  \{W, \lim^V G \circ F^*\} \;\cong\; \lim^V \{W, G \circ F^*\}
\end{equation*}
whenever the relevant weighted limits exist. Symbolically: limits commute
with limits (here weighted limits with weighted limits in the diagram).
\end{proposition}

\begin{proof}
Apply the defining isomorphism, push the hom-functor through the inner
limit (right adjoints preserve limits), and recognize the end-of-end as the
end-over-the-product (Fubini for ends, Chapter~26).
\end{proof}

\begin{remark}
The dual statement — weighted colimits commute with weighted colimits — also
holds, by Fubini for coends. These two statements together make the
\emph{calculus} of weighted (co)limits as smooth as the calculus of ordinary
(co)limits, with one extra dimension of flexibility (the weight).
\end{remark}


% ═══════════════════════════════════════════════════════════════════════════
\section{Existence and Completeness}
\label{sec:wlim-completeness}
% ═══════════════════════════════════════════════════════════════════════════

\subsection{Formalizing}

\begin{definition}[Weighted completeness]
\label{def:wcomplete}
A $\mathcal{V}$-category $\mathcal{C}$ is \term{$\mathcal{V}$-complete} if
$\{W, F\}$ exists for every small diagram $F : \mathcal{J} \to \mathcal{C}$
and weight $W : \mathcal{J} \to \mathcal{V}$. Dually, \term{$\mathcal{V}$-cocomplete}
for weighted colimits.
\end{definition}

\begin{theorem}[Weighted completeness from cotensors and conical limits]
\label{thm:wcomplete-from-pieces}
A tensored and cotensored $\mathcal{V}$-category $\mathcal{C}$ is
$\mathcal{V}$-complete iff its underlying ordinary category is complete (has
all small conical limits). In that case,
\begin{equation*}
  \{W, F\} \;\cong\; \lim_{j \in \mathcal{J}} W(j) \pitchfork F(j),
\end{equation*}
where the outer $\lim$ is the ordinary conical limit and $W(j) \pitchfork F(j)$
is the cotensor.
\end{theorem}

\begin{prooflog}
\proofstep{Compute $\mathcal{C}(c, \{W, F\}) \cong \int_j \mathcal{V}(W j,
\mathcal{C}(c, F j))$.}{End formula.}
\proofstep{$\int_j \mathcal{V}(W j, \mathcal{C}(c, F j)) \cong \int_j
\mathcal{C}(c, W(j) \pitchfork F(j))$.}{By the cotensor adjunction
$\mathcal{V}(V, \mathcal{C}(c, X)) \cong \mathcal{C}(c, V \pitchfork X)$.}
\proofstep{$\int_j \mathcal{C}(c, W(j) \pitchfork F(j)) \cong
\mathcal{C}\!\left(c, \lim_j W(j) \pitchfork F(j)\right)$.}{Hom-functor is
continuous; takes ends (limits) in to limits.}
\proofstep{Therefore $\{W, F\} \cong \lim_j W(j) \pitchfork F(j)$ by Yoneda.}{
Two objects with isomorphic hom-out are isomorphic.}
\proofstep{Existence: if conical limits and cotensors exist, the right-hand
side does, so does the weighted limit.}{Reduction.}
\proofstep{Conversely, taking $W = \Delta I$ recovers conical limits; taking
$\mathcal{J}$ discrete recovers cotensors.}{The two pieces are special
weighted limits.}
\end{prooflog}

\begin{corollary}
A symmetric monoidal closed complete category $\mathcal{V}$ is $\mathcal{V}$-complete
when self-enriched. (Cotensors are internal homs.)
\end{corollary}

\subsection{Reorganizing: Cocompleteness Dually}

\begin{theorem}[Weighted cocompleteness from tensors and conical colimits]
\label{thm:wcocomplete}
A tensored $\mathcal{V}$-category $\mathcal{C}$ with all conical colimits is
$\mathcal{V}$-cocomplete, and
\begin{equation*}
  W \star F \;\cong\; \operatorname{colim}_j W(j) \otimes F(j).
\end{equation*}
\end{theorem}

\begin{proof}
Dual to Theorem~\ref{thm:wcomplete-from-pieces}. Replace ends by coends,
products by coproducts, $\pitchfork$ by $\otimes$, $\lim$ by $\operatorname{colim}$.
\end{proof}


% ═══════════════════════════════════════════════════════════════════════════
\section{Weighted Kan Extensions: A Preview}
\label{sec:wkan-preview}
% ═══════════════════════════════════════════════════════════════════════════

The weighted-limit framework gives an immediate definition of Kan extensions
in the enriched setting. We state the formula here; Chapter~31 develops the
general theory.

\begin{definition}[Weighted Kan extension]
\label{def:wkan}
For a $\mathcal{V}$-functor $K : \mathcal{C} \to \mathcal{D}$ and $F :
\mathcal{C} \to \mathcal{E}$, the \term{(pointwise) left Kan extension}
$\operatorname{Lan}_K F : \mathcal{D} \to \mathcal{E}$ is defined by the
weighted colimit
\begin{equation*}
  (\operatorname{Lan}_K F)(d) \;=\; \mathcal{D}(K-, d) \star F.
\end{equation*}
Dually, the right Kan extension is the weighted limit
\begin{equation*}
  (\operatorname{Ran}_K F)(d) \;=\; \{\mathcal{D}(d, K-), F\}.
\end{equation*}
\end{definition}

\begin{remark}[Coend form]
Using the coend formula:
\begin{equation*}
  (\operatorname{Lan}_K F)(d) \;=\; \int^c \mathcal{D}(K c, d) \otimes F(c),
  \qquad
  (\operatorname{Ran}_K F)(d) \;=\; \int_c [\mathcal{D}(d, K c), F(c)],
\end{equation*}
the latter using cotensor $[-,-]$. These recover the formulas of Chapter~18
(in the ordinary case) and Chapter~24's promised enriched versions.
\end{remark}

\begin{remark}[Every weighted (co)limit is a Kan extension]
A weighted limit $\{W, F\}$ for $F : \mathcal{J} \to \mathcal{C}$ and $W :
\mathcal{J} \to \mathcal{V}$ is the right Kan extension of $F$ along
$W : \mathcal{J} \to \mathcal{V}$ evaluated at $I \in \mathcal{V}$:
\begin{equation*}
  \{W, F\} \;\cong\; (\operatorname{Ran}_W F)(I).
\end{equation*}
Dually for colimits. This identifies (weighted) limit theory with Kan
extension theory — they are the same thing presented from different
viewpoints.
\end{remark}


% ═══════════════════════════════════════════════════════════════════════════
\section{Exercises}
\label{sec:wlim-exercises}
% ═══════════════════════════════════════════════════════════════════════════

\begin{exercisesbox}
\begin{qa}
\qaitem{Define the weighted limit $\{W, F\}$ of $F : \mathcal{J} \to \mathcal{C}$ by $W : \mathcal{J} \to \mathcal{V}$.}{The object with $\mathcal{C}(c, \{W, F\}) \cong [\mathcal{J}, \mathcal{V}]_\mathcal{V}(W, \mathcal{C}(c, F-))$ in $\mathcal{V}$, naturally in $c$.}
\qaitem{Write the end formula for the weighted limit.}{$\mathcal{C}(c, \{W, F\}) \cong \int_j \mathcal{V}(W j, \mathcal{C}(c, F j))$.}
\qaitem{What is the conical limit as a weighted limit?}{$\lim F = \{\Delta I, F\}$ (or $\{\Delta\{*\}, F\}$ in the $\mathbf{Set}$-case): weighted by the constant unit.}
\qaitem{What is the cotensor as a weighted limit?}{For $\mathcal{J} = \mathbf{1}$, $F$ a single object $X$, $W$ a single object $V$: $V \pitchfork X = \{V, X\}$, satisfying $\mathcal{C}(c, V \pitchfork X) \cong \mathcal{V}(V, \mathcal{C}(c, X))$.}
\qaitem{Define the weighted colimit $W \star F$.}{For $W : \mathcal{J}^{\mathrm{op}} \to \mathcal{V}$, the object with $\mathcal{C}(W \star F, c) \cong [\mathcal{J}^{\mathrm{op}}, \mathcal{V}]_\mathcal{V}(W, \mathcal{C}(F-, c))$ in $\mathcal{V}$.}
\qaitem{Coend formula for the weighted colimit?}{$W \star F \cong \int^j W(j) \otimes F(j)$, in a tensored category.}
\qaitem{Conical colimit as a weighted colimit?}{$\operatorname{colim} F = \Delta I \star F$.}
\qaitem{What is the tensor $V \otimes X$ as a weighted colimit?}{For $\mathcal{J} = \mathbf{1}$, the colimit weighted by $V$ of the diagram at $X$. Satisfies $\mathcal{C}(V \otimes X, c) \cong \mathcal{V}(V, \mathcal{C}(X, c))$.}
\qaitem{State the Yoneda reduction for weighted limits.}{$\{\mathcal{J}(j_0, -), F\} \cong F(j_0)$: representable weights pick out values.}
\qaitem{Coend Yoneda reduction for weighted colimits?}{$\mathcal{J}(-, j_0) \star F \cong F(j_0)$.}
\qaitem{What is a comma object as a weighted limit?}{The lax limit of a 2-cell, with weight given by a specific $\mathbf{Cat}$-valued functor $\mathbf{2} \to \mathbf{Cat}$.}
\qaitem{When is a $\mathcal{V}$-category $\mathcal{V}$-complete?}{When $\{W, F\}$ exists for every small diagram $F$ and weight $W$.}
\qaitem{How is weighted completeness reduced to cotensors plus conical limits?}{$\{W, F\} \cong \lim_j W(j) \pitchfork F(j)$ in a tensored/cotensored $\mathcal{V}$-category.}
\qaitem{Dually for weighted cocompleteness?}{$W \star F \cong \operatorname{colim}_j W(j) \otimes F(j)$.}
\qaitem{Why do weighted limits commute with weighted limits (limits commute with limits)?}{Because they are both ends, and ends commute by Fubini.}
\qaitem{Why does the coend formula for weighted colimits require $\mathcal{C}$ to be tensored?}{Because the formula uses $W(j) \otimes F(j)$; without tensors, this is not a well-defined object of $\mathcal{C}$.}
\qaitem{Define the (pointwise) left Kan extension in the weighted setting.}{$(\operatorname{Lan}_K F)(d) = \mathcal{D}(K-, d) \star F = \int^c \mathcal{D}(K c, d) \otimes F(c)$.}
\qaitem{Right Kan extension in the weighted setting?}{$(\operatorname{Ran}_K F)(d) = \{\mathcal{D}(d, K-), F\} = \int_c [\mathcal{D}(d, K c), F(c)]$.}
\qaitem{Is every weighted (co)limit a Kan extension?}{Yes: $\{W, F\} \cong (\operatorname{Ran}_W F)(I)$ and $W \star F \cong (\operatorname{Lan}_W F)(I)$, where $W$ is treated as a functor $\mathcal{J} \to \mathcal{V}$ and the Kan extensions are along it.}
\qaitem{Why do representable weights ``evaluate'' the diagram?}{Because $\mathcal{J}(j_0, -)$ acts as the ``point at $j_0$'' under the Yoneda lemma; the weighted limit reduces by Yoneda to $F(j_0)$.}
\qaitem{What is a homotopy limit in this framework?}{A weighted limit where the weight encodes higher coherence (e.g., simplicial structure). When $\mathcal{V} = \mathbf{sSet}$ and $W$ is a fibrant replacement of the constant simplicial set, $\{W, F\}$ is the homotopy limit.}
\qaitem{What is the relationship between the end of a $\mathcal{V}$-functor and a weighted limit?}{An end is the weighted limit by the hom-profunctor: $\int_c T(c, c) = \{\mathcal{C}(-,-), T\}$ where $T : \mathcal{C}^{\mathrm{op}} \otimes \mathcal{C} \to \mathcal{V}$.}
\qaitem{In $\mathbf{Set}$-enriched setting, what does ``tensored'' mean?}{$\mathbf{Set}$-tensored means having copowers: $S \cdot X = \coprod_S X$ exists. Every cocomplete category is $\mathbf{Set}$-tensored.}
\qaitem{Why is the dual variance condition on $W$ for colimits ($\mathcal{J}^{\mathrm{op}}$ vs.\ $\mathcal{J}$) needed?}{Because the weighted colimit's defining isomorphism contains $\mathcal{C}(F-, c)$ which is contravariant in the diagram source; the weight must match this contravariance.}
\qaitem{What is the underlying ordinary functor of $\{-, -\} : [\mathcal{J}, \mathcal{V}]^{\mathrm{op}} \otimes [\mathcal{J}, \mathcal{C}] \to \mathcal{C}$?}{The ordinary bifunctor taking $(W, F)$ to the weighted limit $\{W, F\}$.}
\qaitem{If $\mathcal{V}$ is cartesian closed, what is the cotensor in $[\mathcal{C}, \mathcal{V}]$?}{$(V \pitchfork F)(c) = V \pitchfork F(c) = [V, F(c)]_\mathcal{V}$, computed pointwise.}
\qaitem{What is a ``descent object'' as a weighted limit?}{A weighted limit of a coherence diagram (a cosimplicial object in $\mathcal{C}$, truncated). It encodes ``descent data'' from a cover to a base.}
\qaitem{Why does flexible weighted limits theory work for 2-categories?}{Because $\mathbf{Cat}$ is closed under (small) products and cotensors with categories; tensors are categories indexed by a category.}
\qaitem{What is the relationship between weighted (co)limits and ordinary natural transformations?}{Weighted limits represent ``natural transformations from the weight to the hom-functor''; the weight becomes a probe for the diagram.}
\qaitem{Why does the chapter's title pair ``weighted'' with ``limits''?}{Because ordinary limits are the case $W = \Delta I$ (a single ``constant'' weight); weighted limits add the flexibility of an arbitrary weight, generalizing the limit notion to all of enriched/2-categorical settings.}
\end{qa}
\end{exercisesbox}


% ═══════════════════════════════════════════════════════════════════════════
\section*{Chapter Summary}
\addcontentsline{toc}{section}{Chapter Summary}
% ═══════════════════════════════════════════════════════════════════════════

\begin{summarybox}
\textbf{Weighted limits via ends.}\quad $\{W, F\}$ is defined by
$\mathcal{C}(c, \{W, F\}) \cong \int_j \mathcal{V}(W j, \mathcal{C}(c, F j))$.
Conical limits are the case $W = \Delta I$.

\textbf{Weighted colimits via coends.}\quad $W \star F = \int^j W(j) \otimes
F(j)$ in a tensored $\mathcal{V}$-category. Conical colimits are the case
$W = \Delta I$. Tensors are the case $\mathcal{J} = \mathbf{1}$.

\textbf{Yoneda reduction.}\quad Representable weights evaluate the diagram:
$\{\mathcal{J}(j_0, -), F\} \cong F(j_0)$ and $\mathcal{J}(-, j_0) \star F
\cong F(j_0)$.

\textbf{Weighted completeness from pieces.}\quad In a tensored/cotensored
$\mathcal{V}$-category, weighted completeness is equivalent to conical
completeness; $\{W, F\} \cong \lim_j W(j) \pitchfork F(j)$.

\textbf{Functoriality and Fubini.}\quad Weighted (co)limits are functorial in
both arguments. Limits commute with limits (Fubini for ends); colimits with
colimits (Fubini for coends). This makes weighted-limit calculus mechanical.

\textbf{Kan extension preview.}\quad The (pointwise) left Kan extension is
$(\operatorname{Lan}_K F)(d) = \int^c \mathcal{D}(K c, d) \otimes F(c)$ and
the right is $\int_c [\mathcal{D}(d, K c), F(c)]$. Every weighted (co)limit
is a Kan extension; every Kan extension is a weighted (co)limit. Full
development in Chapter~31.
\end{summarybox}


% ═══════════════════════════════════════════════════════════════════════════
\section*{Formal Restatement}
\addcontentsline{toc}{section}{Formal Restatement}
% ═══════════════════════════════════════════════════════════════════════════

\begin{formalrestatement}
\textbf{Weighted limit.}\quad
\begin{equation*}
  \mathcal{C}(c, \{W, F\}) \;\cong\; \int_j \mathcal{V}(W j, \mathcal{C}(c, F j))
  \quad\text{in } \mathcal{V}, \text{ naturally in } c.
\end{equation*}

\medskip
\textbf{Weighted colimit.}\quad
\begin{equation*}
  W \star F \;\cong\; \int^j W(j) \otimes F(j),
  \qquad
  \mathcal{C}(W \star F, c) \;\cong\; \int_j \mathcal{V}(W j, \mathcal{C}(F j, c)).
\end{equation*}

\medskip
\textbf{Conical (co)limits.}\quad $\lim F = \{\Delta I, F\}$; $\operatorname{colim} F
= \Delta I \star F$.

\medskip
\textbf{Cotensor, tensor.}\quad $V \pitchfork X = \{V, X\}$ for $\mathcal{J} =
\mathbf{1}$; $V \otimes X = V \star X$.

\medskip
\textbf{Yoneda reduction.}\quad $\{\mathcal{J}(j_0, -), F\} \cong F(j_0)$ and
$\mathcal{J}(-, j_0) \star F \cong F(j_0)$.

\medskip
\textbf{Completeness reduction.}\quad In a tensored/cotensored $\mathcal{V}$-category,
$\{W, F\} \cong \lim_j W(j) \pitchfork F(j)$ and $W \star F \cong \operatorname{colim}_j
W(j) \otimes F(j)$. Weighted (co)completeness reduces to conical (co)completeness
plus (co)tensors.

\medskip
\textbf{Kan extension formulas.}\quad
\begin{align*}
  (\operatorname{Lan}_K F)(d) &= \mathcal{D}(K-, d) \star F = \int^c \mathcal{D}(K c, d) \otimes F(c), \\
  (\operatorname{Ran}_K F)(d) &= \{\mathcal{D}(d, K-), F\} = \int_c [\mathcal{D}(d, K c), F(c)].
\end{align*}

\medskip
\textbf{Universal identification.}\quad Every weighted (co)limit is a Kan
extension; every Kan extension is a weighted (co)limit.

\medskip
\noindent\textit{Chapter~28 makes the universality of Yoneda formal in this
weighted/enriched setting: \emph{the Yoneda embedding is the free cocompletion}
in the precise sense of weighted colimit theory.}
\end{formalrestatement}
